\section{The critical bulk trajectory}\label{sec:trajectory}

We determine the gradient singleton to an error summable in the
center calculation.  The starting scale in Section~\ref{sec:entry}
depends on the target size.  We therefore first control each finite
history by physical torus observables, and then compare histories
where their absolute scales overlap.

Fix the block ratio \(L\), the activity weights, and the regulator
reserves as in Section~\ref{sec:entry}.  Choose
the entry parameter \(R\) sufficiently large, and put \(P=R^5\).
For a future center observable \(M_{L^N}\), use one entry for all
integer targets \(2^k\le N<2^{k+1}\).  For each sufficiently large
\(k\), choose a permitted cell side \(s_k\asymp 2^k\) such that
\begin{equation}\label{trj:starts}
 m_k=qRs_k=L^{j_0(k)},\qquad j_0(k)=O(k).
\end{equation}
The constant in \(\asymp\) may depend on the fixed parameters.  The
divisibility in the entry construction permits this choice: consecutive
permitted sides have the fixed ratio \(L\).  The cell and entry sides
\(s_k,m_k\asymp N\) are lattice lengths; their logarithmic entry index
is only \(j_0(k)=O(\log N)\).  Coordinates with index \(j\) live at
side \(L^j\), and step \(j\) sends them to side \(L^{j+1}\).  Thus
\(j\) is an absolute index, never reset at \(j_0(k)\).  The useful
running indices have \(j\asymp N\); the center calculation stops at
\(j'=N-O(\log N)\), after order \(N\) steps from the entry.

Write \(\mathcal H^{(k)}\) for the plane history obtained by placing the
untwisted entry activity wholly in the remainder and applying the weak
map.  Its unbalanced coordinates are
\((t_j^{(k)},z_j^{(k)},\mathcal R_j^{(k)})\): the two singleton terms are
\(t_j^{(k)}e_B^0/2+z_j^{(k)}c_B^\alpha\), with the normalization in
\eqref{eq:singletons}.  We abbreviate the norm
\eqref{eq:activity-norm} at side \(L^j\), using the observation regulator
for \(s_k\), by \(\|\cdot\|_{j,k}\).  All constants below are uniform
over these norms.  Norms from two different histories will not be
subtracted from each other.

\begin{proposition}[Critical trajectory]\label{trj:trajectory}
The entry choices can be made so that, for every fixed
\(0<c<c'<3\), uniformly for integers
\(c2^k\le j\le c'2^k\) as \(k\to\infty\),
\begin{equation}\label{trj:unbalanced-conclusion}
 \begin{split}
 t_j^{(k)}&=\frac{1}{2(\log L)j}
                  +O\!\left(\frac{\log(2+j)}{j^2}\right),\\
 z_j^{(k)}&=O(j^{-1}),\qquad
 \|\mathcal R_j^{(k)}\|_{j,k}=O(j^{-2}).
 \end{split}
\end{equation}
The constants may depend on \(c,c'\) and the fixed map parameters.
In particular the assertion applies at the stop \(j'=N-O(\log N)\)
in the history chosen for \(2^k\le N<2^{k+1}\).
\end{proposition}

\subsection{Analytic coordinates and a coefficient estimate}

We state precisely the analytic part of the companion recursion used
here.  Within one history omit its superscript \(k\).  Write \(\Phi_j\)
for the complete normalized unbalanced plane weak map on the additive
coordinates above.  It includes the scalar inverse decorations in the
partition algebra of \cite[Proposition~4.8]{B}, and \(\Phi_j(0)=0\).
At the fixed reference \(a=8\pi\), its linearization is
\begin{equation}\label{trj:triangular}
 (t,z,\mathcal R)\longmapsto
 \bigl(t+P_j^t\mathcal R,\ \lambda_jz+P_j^z\mathcal R,
                              \ \mathsf C_j\mathcal R\bigr),
 \qquad \|\mathsf C_j\|\le\vartheta_0<1.
\end{equation}
The projections are bounded uniformly in \(j,k\); the omitted part is
analytic with norm at most \(C(|t|+|z|+\|\mathcal R\|)^2\), with its
Taylor derivative bounds.  These statements are
\cite[Proposition~8.3]{B}, applied with the regulator of
\cite[Lemma~9.2]{B}.  They concern activities at a fixed frequency.
For a compatible torus the singleton coefficients copy the plane
coefficients up to its stop, while its own remainder satisfies the
same contraction and analytic bounds.  No equality between the plane
and torus remainders is part of this assertion.

For our polynomial shells, \(\lambda_j=L^2
e^{-a\Gamma_{L^j}(0,0)/2}\).  The continuum value is one.  The
fixed nearest-neighbor estimate in \cite[Lemma~8.1 and its proof]{B},
with the independent microscopic piece omitted as in
\cite[Equation~(9.11)]{B}, gives
\begin{equation}\label{trj:linear-error}
 |\lambda_j-1|\le CL^{-c_1j},\qquad
 \ell_j:=\prod_{h\ge j}\lambda_h=1+O(L^{-c_1j})
\end{equation}
after increasing the first absolute scale.  Here and below an exponent
may be decreased without changing notation.  Put
\(\mathsf C_{h:j}=\mathsf C_{h-1}\cdots\mathsf C_j\) for \(h>j\) and
\(\mathsf C_{j:j}=I\).  The balanced coordinates are exactly
\begin{equation}\label{trj:balance}
 T=t+\sum_{h\ge j}P_h^t\mathsf C_{h:j}\mathcal R,
 \qquad
 Z=\ell_jz+\sum_{h\ge j}\ell_{h+1}P_h^z\mathsf C_{h:j}\mathcal R.
\end{equation}
The series have geometric operator tails.  These are
\cite[Equations~(8.17)--(8.18)]{B} at the marginal eigenvalue one.
In particular the small linear errors in \(\lambda_j\) have already
been balanced; they are not included in a nonlinear activity error.
The formulas for the maps and the shells above a common absolute
side are independent of the entry side.  The entry-dependent regulator
is used only to bound these formulas.

We use one neutral terminal response both for an exact Gaussian
trajectory and for the physical histories below.  Fix a sufficiently
large dyadic \(D_0\), exceeding all causal padding radii, and set
\begin{equation}\label{trj:stopping-tori}
                       n_j=D_0L^j.
\end{equation}
For each trajectory used below, a torus stopped at side \(L^j\) is
initialized separately at its entry side on the compatible grid of
side \(n_j\).  It has \(D_0^2\) terminal blocks, and the fixed buffer
makes its singleton coefficients copy the plane coefficients through
that stop.
This construction uses only the local map.  The physical cutoff tori
introduced below will also have to satisfy \(n_j\ge s_kP\).

On this torus let \(A\) be the unit-conductance Laplacian, and let
\(\E_j^{\rm ter}\) integrate a mean-free Gaussian \(\psi\) with covariance
\(C_{\ge L^j}=A^+\mathsf P_{L^j}\) and an independent constant
\(\theta\), uniform modulo \(\omega\); write \(\phi=\psi+\theta\).
Let \(q_n(x)=\cos(2\pi x_1/n)\), with the origin fixed by the compatible
grid, and put \(F_{n,\gamma}(w)=w\gamma q_n\).  Here \(\gamma\) ranges
over a fixed compact subinterval of \((0,\infty)\).  Since
\(n=n_j=D_0L^j\), the field norms of \eqref{eq:point-norm} satisfy
\(\sup_{\gamma,X}\|\gamma q_{n_j}\|_{L^j,D_{L^j}(X)}\le C\).
Choose a fixed complex source disk \(|w|<w_{\rm ter}\) with
\(w_{\rm ter}C<h_0\) and a fixed strict margin.  All energy pairings
involving complex sources below are the complex bilinear extensions
of the real lattice pairing.  Thus
\(\ip{F_{n,\gamma}(w)}{AF_{n,\gamma}(w)}
 =w^2\gamma^2\ip{q_n}{Aq_n}\).

At this fixed grid, and for each entry regulator separately, let \(K\)
range over complex local \(\omega\)-periodic analytic activities whose
field dependence has the support collars in Proposition~\ref{prop:local-map}.
Write its unbalanced coordinates as \((t,z,\mathcal R^{\rm tor})\),
and use the open complex ball
\(\nu=|t|+|z|+\|\mathcal R^{\rm tor}\|<\rho_{\rm ter}\), with the
current torus norm for that entry.  All activity variations below are
taken within one such space.  The finite subset sum and the source
Taylor series controlled by \(h_0\) make
\(\E_j^{\rm ter}\cZ_{L^j}(K;\phi+F_{n_j,\gamma}(w))\) jointly analytic
in the activity coordinates and \(w\).  Its absolute expansion on the
fixed \(D_0^2\)-block grid bounds its difference from one by
\(CD_0^2\nu\), uniformly on the source disk.  Choose
\(CD_0^2\rho_{\rm ter}<1/2\).  Both the shifted and unshifted integrals
are then nonzero, so define
\begin{equation}\label{trj:neutral-response}
 \mathcal W_j(K;F)=
 \log\frac{\E_j^{\rm ter}\cZ_{L^j}(K;\phi+F)}
               {\E_j^{\rm ter}\cZ_{L^j}(K;\phi)}.
\end{equation}
The common phase has the same normalized law in both integrals, and
the logarithm takes the branch equal to zero at \(F=0\).  On fixed
smaller concentric activity and source domains, joint Cauchy bounds
control every fixed input and source derivative uniformly in the
entry and stop.  The tubes below lie in such an inner activity ball;
its convexity keeps the segments between actual and pure states inside
the larger analytic domain.

For the half-period action
\((\mathcal S K)(X,\phi)=K(X,\phi+\omega/2)\), applied to every
constituent together, the activity ball is invariant.  Changing the
common phase in both integrals gives
\(\mathcal W_j(\mathcal S K;F)=\mathcal W_j(K;F)\).
On the inner domains the terminal expansion gives, with two source
derivatives,
\begin{equation}\label{trj:neutral-readout}
 \begin{split}
 \mathcal W_j(K;F_{n_j,\gamma}(w))
   &=\frac t2\ip{F_{n_j,\gamma}(w)}{AF_{n_j,\gamma}(w)}
             +O(\|\mathcal R^{\rm tor}\|+\nu^2),\\
 \left.\partial_w^2\mathcal W_j(K;F_{n_j,\gamma}(w))\right|_{w=0}
   &=\kappa_{n_j,\gamma}t+O(\|\mathcal R^{\rm tor}\|+\nu^2),\\
 \kappa_{n,\gamma}
   &=\gamma^2\ip{q_n}{Aq_n}
     =2\gamma^2n^2\sin^2(\pi/n)\asymp1.
 \end{split}
\end{equation}
In particular this reads the actual unbalanced gradient \(t\).
To obtain the expansion, retain the terms linear in \(K\) in the
terminal subset sum.  The centered Gaussian removes the cross term in
the gradient energy, and the common phase removes the fundamental
singleton even after the shift.  The linear remainder costs its norm;
the remaining terms cost \(O(\nu^2)\) by the finite block count and
the regulator reserve.  This is the calculation in
\cite[Equations~(10.2)--(10.3)]{B}, here for the specified first cosine.

For each chosen entry side there is also an exact Gaussian trajectory.
At side \(m\), for every nonempty polymer \(X\) connected in the
subset convention, define the local occupied-component input by
\[
 g_B(d,\phi)=e^{d e_B^0(\phi)/2}-1,\qquad
 K_m^g(X,\phi;d)=\prod_{B\in X}g_B(d,\phi).
\]
For every occupied set \(V\), the product over its components is
\(\prod_{B\in V}g_B\).  Hence on each finite grid
\(\cZ_m(K_m^g;\phi)=\prod_B(1+g_B)
=\exp\{d\sum_Be_B^0(\phi)/2\}\).  Use the same local input on the
plane and compatible tori.  With the fixed activity weights, a common
sufficiently small \(d\)-disk makes this an allowed analytic input:
its linear term is the singleton \(d e_B^0/2\), and its remaining
singleton and multi-block terms have norm \(O(d^2)\).  Every factor
is invariant under common field translations, so the map preserves
the vanishing of all charged coordinates.  Denote the balanced
gradient, unbalanced gradient, and unbalanced remainder at scale \(j\)
by \(T_j^g(d)\), \(t_j^g(d)\), and \(\mathcal R_j^g(d)\).
On a fixed complex disk,
uniformly in all sufficiently large entry sides and all later steps,
\begin{equation}\label{trj:gaussian-curve}
 T_j^g(d)=d+O(d^2),\qquad Z_j^g(d)=0,
 \qquad \|\mathcal R_j^g(d)\|\le C|d|^2.
\end{equation}
The bounds are analytic, the functions \(T_j^g\) have inverses on a
common smaller disk, and the step carries the point with parameter
\(d\) to the point with that same parameter.  The fixed-entry curve is
\cite[Lemma~10.3]{B}; we record why its estimates are uniform as the
entry side varies.  For entry side \(m\) and each stop \(j\), let
\(K^{g,\rm tor}_{j|j}(d)\) be the terminal activity of the separately
initialized pure torus.  Apply the exact local-map identities to
\(\phi+F\), where \(F=F_{n_j,\gamma}(w)\).  The field-independent
extraction factors cancel in \eqref{trj:neutral-response}.  The shells
from \(m\) to \(L^j\) and the terminal covariance telescope to
\(C_{\ge m}\) on mean-free modes, and
\(\sum_B e_B^0(\phi)=\ip{\phi}{A\phi}\) on the torus.  Gaussian
completion thus gives, while the terminal state lies in the readout
ball, the exact identity
\begin{equation}\label{trj:pure-response}
 \mathcal W_j(K^{g,\rm tor}_{j|j}(d);F)
 =\frac d2\ip{A^{1/2}F}{
       (I-dA^{1/2}C_{\ge m}A^{1/2})^{-1}A^{1/2}F}.
\end{equation}
The gradient energy and the mean-free shift are insensitive to the
common constant.
On mean-free modes \(A^{1/2}C_{\ge m}A^{1/2}=\mathsf P_m\) lies
between zero and \(I\).  For \(|d|<\delta<1\) the Gaussian determinant
is nonzero and cancels in the normalized ratio, while the resolvent
has norm at most \((1-\delta)^{-1}\).  The second source derivative
of \eqref{trj:pure-response} is consequently \(O(|d|)\), uniformly in
the entry and stop.  These pure tori are available at every stop after
their entry; they do not use the physical cutoff sewing.

Before a first exit from a ball of radius \(\rho\), the bounded map
gives every candidate plane and stopped-torus state full size at most
\(C\rho\).  Choose \(\rho\) so that these states lie inside the fixed
inner readout ball.  Remainder contraction bounds both the plane and
stopped-torus remainders
by \(C(|d|^2+\rho^2)\).  In this pure case the fundamental coordinate
vanishes.  Comparing the second derivative in
\eqref{trj:pure-response} with \eqref{trj:neutral-readout}, whose
multiplier is bounded below and whose scalar is the plane \(t_j^g\)
by copying, gives
\[
 |t_j^g|\le C\bigl(|d|+\|\mathcal R^{g,\rm tor}_{j|j}\|
                +( |t_j^g|+\|\mathcal R^{g,\rm tor}_{j|j}\|)^2\bigr)
       \le C(|d|+\rho^2).
\]
The candidate full-size bound makes the squared term \(O(\rho^2)\).
The constants are independent of the entry and stage.  Choose
\(\rho\) small and then one disk radius
\(\delta>0\) so that these bounds exclude every first exit for
\(|d|\le\delta\).  Every finite composition is consequently analytic
and uniformly bounded on that common disk.  The triangular
linearization gives
\(\mathcal R_j^g(0)=(\mathcal R_j^g)'(0)=0\) and
\(T_j^g(0)=0,(T_j^g)'(0)=1\).  Cauchy estimates on a smaller common
disk give \eqref{trj:gaussian-curve} and its derivative bounds;
one further shrink gives the common inverse disk.  This is the
no-exit mechanism in \cite[Section~5.4, proof of Proposition~5.1]{C},
using here only the neutral contraction at the marginal parameter.

Define, for this history,
\begin{equation}\label{trj:normal-coordinates}
 u_j=(T_j^g)^{-1}(T_j),\qquad
 Q_j=\mathcal R_j-\mathcal R_j^g(u_j),\qquad
 Q_j=Q_{e,j}+Q_{o,j},
\end{equation}
where \(Q_e,Q_o\) are the even and odd parts under the common shift
\(\phi\mapsto\phi+\pi/\alpha\).  The pure Gaussian points in these
coordinates are \((u,0,0)\) with \(u\) unchanged by the step.  Write
\(Y_j=|Z_j|\).

\begin{lemma}[Analytic parity interface]\label{trj:parity-interface}
There are a fixed small tube, \(\vartheta<1\), and \(C<\infty\),
uniform over the entry sides above, on which the exact map satisfies
\begin{align}
 u_+-u={}&-H_jZ^2+
 O\bigl((|u|+Z^2)(Z^2+\|Q_e\|)+\|Q_e\|^2
                    +\|Q_o\|(|Z|+\|Q_o\|)\bigr),
                                      \label{trj:u-map}\\
 Z_+-Z={}&-b_juZ+
 O\bigl((u^2+Z^2+\|Q_e\|)|Z|
       +( |u|+Z^2+\|Q_e\|+\|Q_o\|^2)\|Q_o\|\bigr),
                                      \label{trj:z-map}\\
 \|Q_{e,+}\|\le{}&\vartheta\|Q_e\|
                         +C(|Z|+\|Q_o\|)^2,
                                      \label{trj:even-map}\\
 \|Q_{o,+}\|\le{}&\vartheta\|Q_o\|
                         +C(|u|+Z^2+\|Q_e\|)|Z|.
                                      \label{trj:odd-map}
\end{align}
The error terms have the differentiated Taylor bounds of the
displayed orders.  The coefficients, independent of the entry side
at a common absolute \(j\), obey fixed positive upper and lower
bounds for sufficiently large \(j\).  Their limits are \(H>0\) and
\begin{equation}\label{trj:b-value}
                       b_* =\frac{a}{4\pi}\log L=2\log L.
\end{equation}
\end{lemma}

\begin{proof}
This is the analytic consequence of \cite[Lemma~10.6]{B} used here.
That lemma is formulated for the companion's true-reference histories;
its proof uses only \eqref{trj:triangular}, the bounded changes
\eqref{trj:balance}--\eqref{trj:normal-coordinates}, and half-period
symmetry.  All those hypotheses hold for the entry activities of
Section~\ref{sec:entry}.  The diagonal balanced linearization removes
linear forcing of \(Z\) into the remainder.  The exact Gaussian curve
removes every nonlinear term that would remain with \(Z=Q=0\).
Parity makes the first coordinate even and the second odd; it gives
the two displayed quadratic terms and the indicated orders for the
two remainder parts.  Terms linear in a remainder with a small
coefficient are absorbed in \(\vartheta<1\).  This recovers precisely
\eqref{trj:u-map}--\eqref{trj:odd-map} as estimates on the analytic
activity map.  More specifically, \(u=T+O(T^2)\) depends only on \(T\),
\(Q=\mathcal R+O(u^2)\), and \(Z\) is unchanged, while the balanced
scalar linear rows contain no remainder.  These substitutions can
alter a pure quadratic gradient term but cannot alter the \(Z^2\) and
\(uZ\) coefficients.  Thus
\(H_j=-b_j^{zz}\) and \(b_j=-b_j^{tz}\) are the balanced map coefficients
of \cite[Proposition~8.5]{B}; their signs and limits, and
\eqref{trj:b-value}, follow there and in \cite[Equation~(8.22)]{B}.
\end{proof}

For the precision in Proposition~\ref{trj:trajectory} we need a rate
for these two map coefficients.  The next estimate is a refinement
of the finite-depth argument in the proof of the cited proposition.

\begin{lemma}[Rate for the quadratic coefficients]\label{trj:coefficient-rate}
For the fixed nearest-neighbor reference and fixed \(L\), there is
\(C<\infty\) such that, at every sufficiently large absolute index,
\begin{equation}\label{trj:coefficient-bound}
                 |H_j-H|+|b_j-b_*|\le C(1+j)^{-2}.
\end{equation}
\end{lemma}

\begin{proof}
A quadratic coefficient is the Hessian of the normalized unbalanced
map followed by the linear balancing row at the output.  We make the
truncation of that row explicit.  Let
\(\pi_t,\pi_z,\pi_{\mathcal R}\) be the coordinate projections, and let
\(\mathbf t_j=(1,0,0)\), \(\mathbf z_j=(0,\ell_j^{-1},0)\) be the
unbalanced inputs corresponding to unit balanced scalar directions at
zero remainder.  All derivatives of \(\Phi_j\) here are at zero in
the additive weak coordinates \(t e_B^0/2+z c_B^\alpha+\mathcal R\).
They include the regrouping by the scalar decorations
\(d_B=(1+a_B)^{-1}-1\) already inside the normalized map; the external
extracted product is not an additional
coordinate row.  The exponential initialization above is used for the
separate Gaussian curve.  The operators \(P_j^t,P_j^z,\mathsf C_j\)
are the three components of \(D_{\mathcal R}\Phi_j(0)\) in
\eqref{trj:triangular}.

For an absolute output index \(i\) and \(h\ge0\), set
\begin{equation}\label{trj:truncated-balance}
 \mathsf S^t_{i,h}=\sum_{r=0}^{h-1}P^t_{i+r}\mathsf C_{i+r:i},
 \qquad
 \mathsf S^z_{i,h}=\sum_{r=0}^{h-1}
            \ell_{i+r+1}P^z_{i+r}\mathsf C_{i+r:i},
\end{equation}
with empty sums zero.  The exact truncated coefficients are
\begin{equation}\label{trj:truncated-coefficients}
 \begin{split}
 H_{j,h}
  &=-\tfrac12(\pi_t+\mathsf S^t_{j+1,h}\pi_{\mathcal R})
       D^2\Phi_j(0)[\mathbf z_j,\mathbf z_j],\\
 b_{j,h}
  &=-(\ell_{j+1}\pi_z+\mathsf S^z_{j+1,h}\pi_{\mathcal R})
       D^2\Phi_j(0)[\mathbf t_j,\mathbf z_j].
 \end{split}
\end{equation}
Here \(D^2\Phi_j(0)\) is the bilinear Hessian: the \(Z^2\) Taylor
coefficient has the factor \(1/2\), whereas the mixed \(TZ\) coefficient
does not.  The chain rule for the two linear balancing maps, followed
by the coefficient identification in Lemma~\ref{trj:parity-interface},
gives these formulas from \cite[Proposition~8.5]{B}.  At \(h=0\) they
retain the direct Hessian row; for \(h>0\) the last future projection
is \(P_{j+h}\).  Thus the chain of local \(\mathsf C\) and \(P\)
operations uses shell indices \(j,\ldots,j+h\).  Each \(\ell\) is
retained as one scalar factor; its infinite product is controlled
separately by \eqref{trj:linear-error}.

Let \(q_{j,h}\) denote either coefficient in
\eqref{trj:truncated-coefficients}, with the same choice throughout.
The geometric tail of the output row gives
\(|q_j-q_{j,h}|\le C\vartheta_0^h\).  For fixed \(h\), the finite-depth
convergence in the proof of \cite[Proposition~8.5]{B} gives
\(q_{\infty,h}:=\lim_{j\to\infty}q_{j,h}\).  Its termwise continuum
evaluation is the comparison used below.  Taking the fixed-depth limit
in the same geometric tail bound gives
\(|q_\infty-q_{\infty,h}|\le C\vartheta_0^h\), where \(q_\infty\) is
\(H\) or \(b_*\).  This definition uses no continuum activity map.

Evaluate the shells separately along this chain, so
only the fixed scale ratio \(L\) occurs; no kernel estimate at the
growing ratio \(L^h\) is needed.  For a shell at side \(L^{j+r}\),
\(0\le r\le h\), use scaled jets
\((L^{j+r})^d\nabla^d\Gamma_{L^{j+r}}\).  The Fourier proof of
\cite[Lemma~8.1]{B}, with the factor for the independent microscopic
piece omitted, gives a uniform bound for these jets and an error
\(CL^{-c_2j}\) against their continuum jets, through order eight.
The error is uniform in the spatial arguments: the proof bounds the
absolute Fourier integral after splitting at a small power of the
side.  The Wick entries below use total difference order at most two,
so their one further spatial derivative is among these bounded orders.
The same proof gives \eqref{trj:linear-error}; hence
replacing any balance multiplier \(\ell_{j+r}\) or its inverse by one costs
\(CL^{-c_2j}\).

Here is a bound for the complexity of the truncated expression.
At Taylor order two in the activity inputs, the exact algebra
\cite[Proposition~4.8]{B} uses only products of at most two singleton
inputs and the degree at most two terms of each scalar inverse.
The resulting supports have a fixed maximum number of blocks.  Every
subsequent propagation in the projection series is linear in that
quadratic image.  A residual keeps its coarse closure and a localized
piece is a singleton; \(|\overline X|\le |X|\).  The maximum block
cardinality therefore remains fixed.  At each step, for a prescribed
output root, there are at most \(C_L\) possible shapes, placements,
and choices of linear localization.  This accounts for the block
branches even when all pieces in a subtraction are kept separately.

The field dependence before the first such projection is a polynomial
in gradients of degree at most four times at most two point characters.
Neutral localization evaluates at zero and takes at most two affine
derivatives; fundamental localization evaluates on a constant and
keeps a character.  Thus propagation either keeps a dependence of
this form or replaces it by a scalar, a quadratic gradient polynomial,
or a character.  Gaussian expectations introduce only covariances of
these jets and covariances with the characters.  For example, with
\(d\le4\), real Gaussian linear forms \(L_r\zeta\), and a real linear
character argument \(V\),
\[
 \E\!\left[\prod_{r=1}^d L_r\zeta\,e^{iV}\right]
 =e^{-\Var(V)/2}
   \sum_{\pi}
    \prod_{\{r,s\}\in\pi}\Cov(L_r\zeta,L_s\zeta)
    \prod_{r\text{ unpaired}} i\Cov(L_r\zeta,V),
\]
where \(\pi\) ranges over partial pairings.  This formula also covers
constant and affine backgrounds after expansion.  It involves no
inverse of a covariance matrix.  Phase integration retains only the
charges of the at most two original characters.  Each step adds a
bounded number of polynomial or covariance factors to a term already
computed; retaining earlier site sums as coefficient indices, a product
has at most \(C(h+1)\) factors.  Distributing all Wick pairings and
Taylor choices, as well as the block branches, is bounded by
\(\exp\{C_L(h+1)^2\}\).  This generous bound also covers repeated
expansion of the accumulated coefficients.  The variance exponentials
in the displayed formula have modulus at most one, and
\(|e^{-v/2}-e^{-w/2}|\le |v-w|/2\) for \(v,w\ge0\).

All remaining spatial indices are normalized sums on half-open
rectangular blocks.  For the energy singleton write, for example,
\(\sum_x(\nabla\phi_x)^2=L^{-2j}\sum_x(L^j\nabla\phi_x)^2\), so it has
the same area normalization as the character singleton.  Localization
adds sums of that form and bounded block assignments.  Expressing
directions or coordinates from one of the first \(h+1\) scales in
another scale costs at most a fixed power of \(L^h\).  Consequently
each unmultiplied entry, and its rectangle-wise Lipschitz bound in one
scaled spatial variable, is at most \(\exp\{C_L(h+1)\}\).  This follows
from the jet bounds above for covariance entries, and directly for
the affine values.  A fixed neighbor stencil moves a point by one mesh
step; its comparison uses the bound for the normalized difference jet
and one spatial derivative of the limiting jet, rather than a derivative
of the lattice error.  Rectangular Riemann summation therefore costs
one mesh step at a boundary or in a Lipschitz estimate.  There are at
most \(C(h+1)\) successive normalized sums in a term, their total masses
are bounded by the preceding scale factors, and the smallest side is
\(L^j\).

Applying product differences to at most \(C(h+1)\) factors and then
performing these Riemann comparisons one index at a time proves, for
each of the two truncated coefficients,
\begin{equation}\label{trj:finite-depth-error}
 |q_{j,h}-q_{\infty,h}|
       \le L^{-c_3j}\exp\{C_L(h+1)^4\}.
\end{equation}
Indeed the elementary jet or mesh error supplies \(L^{-c_3j}\);
the bounds for products, branches, coordinate changes, and the
finitely many rectangle boundaries just proved are each bounded by
an exponential of a polynomial of degree at most two in \(h+1\).
Their product is bounded by the displayed larger exponent.  Balance
multiplier errors from \eqref{trj:linear-error} satisfy the same bound.

Take \(h=\lceil C_4\log(2+j)\rceil\), with \(C_4\) large enough that
\(\vartheta_0^h\le C(1+j)^{-3}\).  The right side
of \eqref{trj:finite-depth-error} is at most \(C(1+j)^{-3}\) for all
sufficiently large \(j\).  Adding the two geometric tails proves
\eqref{trj:coefficient-bound}, after enlarging \(C\) for finitely many
indices.  This argument uses only the explicit Gaussian kernels and
the finite-order local algebra.  No convergence rate for a height
scaling limit has been used.
\end{proof}

\subsection{Physical constraints from stopping tori}

For history \(k\), the valid physical stops on the compatible cell and
block grids are
\begin{equation}\label{trj:valid-stops}
 \mathcal J_k=\{j\in\Z:j_0(k)\le j\le4\cdot2^k,\ n_j\ge s_kP\}.
\end{equation}
For every \(j\in\mathcal J_k\), initialize the exact cutoff torus of
side \(n_j\) separately at \(m_k\) and run it to side \(L^j\).
At a stage \(j_0(k)\le i\le j\), its scalar coordinates are the plane
values \(t_i^{(k)},z_i^{(k)}\), by copying their padded neighborhoods.
Write \(\mathcal R_{i|j}^{\rm tor,(k)}\) for its own remainder.
Its norm is the torus version of \(\|\cdot\|_{i,k}\) with the same
entry regulator.  We suppress \(k\) while working within one history.
No torus is continued from a different size.  When \(n_j<s_kP\) only
plane steps are needed.
There are at most
\(\max\{0,\lceil\log_L(R^4/(D_0q))\rceil\}\) of them.  The bounded
map amplifies the entry norm \(Ce^{-cR}\) by only \(R^{O(1)}\) on
these steps, so increasing \(R\) controls them.  These are the stopping
conventions and the initial estimate in \cite[Section~10.1]{B}.

Let \(\E_n^{\rm tor}\) denote the original periodic Bessel height law
on \((\Z/n\Z)^2\), with its constant modulo \(2\pi\).  Put
\(f(x)=\cos(2\pi x_1)\) on the unit torus and
\begin{equation}\label{trj:probes}
 \mathfrak g=\ip{f}{(-\Delta_{\T^2})^{-1}f}>0,\qquad
 x_n=\ip{h}{n^{-2}f(\cdot/n)},\qquad
 X_{n,s}=\ip{y}{f_{n,s}},
\end{equation}
where \(f_{n,s}(Q)=(s/n)^2 f(x_Q/n)\).  Here \(x_Q\) is the arithmetic
center of the primary sites of \(Q\), in the same coordinates as
\(q_n\): for \(Q=v_Q+\{0,\ldots,s-1\}^2\),
\(x_Q=v_Q+\tfrac{s-1}{2}(1,1)\).  Both tests have zero total.  Also write
\(\chi_{n,s}=e^{i\bar y}\), where \(\bar y\) is the mean of the cell
observations, and write \(e^{i\bar h}\) for the uniform height character.

\begin{lemma}[Torus height diagnostics for the XY input]\label{trj:physical-torus}
Along the compatible even tori \(n\to\infty\),
\begin{align}
 \E_n^{\rm tor}e^{w x_n}&\longrightarrow
                  e^{a\mathfrak g w^2/2},\label{trj:torus-laplace}\\
 \E_n^{\rm tor}[e^{i\bar h}e^{w x_n}]
   -\E_n^{\rm tor}e^{i\bar h}\,e^{a\mathfrak g w^2/2}
       &\longrightarrow0.\label{trj:torus-factorization}
\end{align}
The first assertion includes convergence of every fixed moment, and
both assertions are locally uniform in complex \(w\).  They also hold
for \(X_{n,s},\chi_{n,s}\) when \(s/n\to0\).  For any prescribed small
error in a fixed finite set of these diagnostics, one can first take
\(P\) sufficiently large and then \(s\) sufficiently large so that
the error is small uniformly over compatible \(n\ge sP\).
\end{lemma}

\begin{proof}
We first prove the mean-free assertion using the XY height inputs.
For a smooth periodic shift \(v\) let
\(X_n^v=\ip{h}{A_{\T_n}v(\cdot/n)}\), the sum of
\(\nabla h\,\nabla v(\cdot/n)\) over primary edges.  Cut the torus
into a fixed array of macroscopic rectangles and keep the internal
parts of this test.  The cross-seam part has \(O(n)\) edges with
coefficients \(O(n^{-1})\), with constants depending on the fixed
array, and hence has reference cost \(O(n^{-1})\).  The reference-cost
bound and centered comparison of \cite[Lemmas~2.2--2.3]{C}, followed
by the free rectangle limits \cite[Theorem~2.4]{C}, give
\[
 \limsup_n\Var_{\rm tor}X_n^v\le a\int_{\T^2}|\nabla v|^2.
\]
One can equivalently subdivide further and approximate the internal
gradients by constants: the additional squared cost is controlled by
their modulus of continuity, exactly as in the partial-test argument
of that theorem.  Thus the inequality is a quadratic-form bound for
every finite list of smooth shifts.

If \(v\) has support compactly contained in an unwrapped rectangle,
add a positive-width primary pinned rim inside that rectangle, separated
from the support and surrounding it.  The test is supported inside the
rim, and the fixed rim separates that part of the law from the rest of
the torus.  Centered pin comparison bounds its original variance from
below by the pinned variance.  The mixed limit \cite[Theorem~3.1]{C} gives
for the latter \(a\int|\nabla v|^2\): in its pinned energy space, the
Riesz representative of the weak Laplacian test is \(v\) itself.
Such shifts therefore attain the upper bound.  On any convergent
finite covariance subsequence, the defect from the upper quadratic
form is positive semidefinite.  A vector with zero defect diagonal
has zero defect pairing with every vector, by Cauchy--Schwarz for that
form.  A partition of unity expresses any smooth torus shift as a finite
sum supported in unwrapped rectangles.  The upper bound is consequently
attained by every smooth shift.

For the Laplace lower bound use the centered inequality
\(\log\E e^X\ge\Var(X)/2\) in \cite[Lemma~2.2]{C}.  For the upper
bound, apply H\"older to the internal and seam parts with a fixed
exponent \(p>1\).  Releasing the seams increases the centered transform
of the internal part and makes the rectangle factors independent.
To apply the free rectangle limits to the internal part, fix a further
macroscopic subdivision and approximate \(\nabla v\) by a constant on
each rectangle.  The squared reference cost of the difference is
bounded by the squared modulus of continuity of \(\nabla v\); the new
seam part again has vanishing cost as \(n\to\infty\).  The free
rectangle and partial-test Laplace argument in \cite[Theorem~2.4]{C}
applies to these constant-gradient tests.  The other H\"older factor
tends to one by the primary reference-cost exponential bound and the
vanishing seam cost.  Take the lattice limit first, then let the
subdivision mesh tend to zero.  The upper limit is at most
\(pa\int|\nabla v|^2/2\); finally let \(p\downarrow1\).  Applying this argument
to real linear combinations and using the uniform fixed-argument
exponential bounds gives Gaussian joint limits and all fixed moments.
The sampled cosine is the weak Laplacian of a smooth periodic cosine
up to a vanishing reference-energy error (or use its exact discrete
eigenvalue \(4\sin^2(\pi/n)\)).  This proves
\eqref{trj:torus-laplace} and its moment assertion.

For completeness, Gaussian convergence of the mean-free test alone
does not imply \eqref{trj:torus-factorization}.  Use the chain
subdivision underlying \cite[Lemma~2.2]{C} on each fixed primary torus,
with pinned representatives.  For two independent fields put
\(S=h_1+h_2\), \(D=h_1-h_2\).  Given their common parity class modulo
\(4\pi\Z\), \(S,D\) are independent with the same centered quadratic
law.  With \(M\) subdivision bonds on each primary edge, write
\(x_{n,M}\) for the same test evaluated on the primary heights and
\(\bar h_{\rm prim}\) for their uniform mean.  Set
\(\mathfrak F_{n,M}(t)=\E e^{i\bar h_{\rm prim}+itx_{n,M}}\).
If \(V\) is the conditional second moment of the test
of \(S\), the replica identity \cite[Equation~(2.12)]{C} gives
\[
 \E V=2\Var(x_{n,M}),\qquad \Var V=2\kappa_4(x_{n,M}).
\]
Differentiation in the two original replicas and conditioning give
\[
 -2\{\mathfrak F_{n,M}\mathfrak F_{n,M}''-(\mathfrak F_{n,M}')^2\}
       =2\Var(x_{n,M})\mathfrak F_{n,M}^2+\mathcal E_{n,M}(t),
\]
where on real \(t\) the error is bounded by
\(\sqrt{2\kappa_4(x_{n,M})}\).  On compact complex sets the same bound has
a fixed exponential-moment factor, by Cauchy--Schwarz.  First pass
through the chain limit at this fixed graph, as in the cited identity.
This gives the analogous bound for \(\mathfrak F_n(t)=\E_n^{\rm tor}
e^{i\bar h+itx_n}\).  Then \eqref{trj:torus-laplace} and fourth-moment convergence show that
every locally uniform entire subsequential limit \(\mathfrak F\) satisfies
\[
 \mathfrak F\mathfrak F''-(\mathfrak F')^2=-a\mathfrak g\mathfrak F^2.
\]
The required entire bounds follow from the primary reference costs.
If \(\mathfrak F\) is not zero identically, solving this equation where
\(\mathfrak F\ne0\) and then continuing analytically shows that
\(\mathfrak F\) is an exponential
quadratic and has no zero.  Translation by half the torus sends the
cosine test to its negative and preserves the mean character; hence
\(\mathfrak F'(0)=0\).  We obtain
\(\mathfrak F(t)=\mathfrak F(0)e^{-a\mathfrak g t^2/2}\).
The identically zero limit obeys the same equality.  Continuation to
\(t=-iw\) proves \eqref{trj:torus-factorization}.  This is the replica
argument of \cite[Lemma~10.1]{B}, now justified for the Bessel law by
the fixed-graph chain passage.

Finally, the coefficient difference between the height part of
\(X_{n,s}\) and \(x_n\) is at most \(Cn^{-2}s/n\) at each site and has
zero total.  Since \(\|A_{\T_n}^+\|_{2\to2}\le Cn^2\), its reference
cost is at most \(C(s/n)^2\).  The observation-noise variance is also
\(O((s/n)^2)\), and its mean character is an independent positive
factor \(e^{-(s/n)^2/2}\).  Reference-cost exponential bounds and
Cauchy--Schwarz control the resulting differences of transforms and
fixed derivatives on compact sets by a quantity tending to zero with
\(s/n\).  This proves the assertions for observations.  Ordinary
convergence in \(n\), with \(s/n\le P^{-1}\), proves the stated order
of uniform choices.  It asserts no rate for that convergence.
\end{proof}

For \(j\in\mathcal J_k\), let \(\widehat\E_{j,k}\) be the normalized
cutoff expectation on the torus \eqref{trj:stopping-tori} in history
\(k\).  The cutoff bound of
Section~\ref{sec:entry}, summed over the faces of these tori, changes
bounded phase observables and the fixed smooth Laplace diagnostics by
at most
\begin{equation}\label{trj:cutoff-error}
 \delta_k\le \exp(-c_5 4^k),\qquad j\le4\cdot2^k,
\end{equation}
for large \(k\).  Indeed their log volume is \(O(2^k)\), while
\(s_k^2\asymp4^k\).  The same statement for the fixed source derivatives
follows from the tilted estimate and its moment bound.

We now identify the physical Laplace correction with the neutral response
\eqref{trj:neutral-response}.  Let \(K_{j|j}^{\rm tor}\) be the terminal
activity at a valid stop, with remainder \(\mathcal R_{j|j}^{\rm tor}\),
and set
\(\nu_j=|t_j|+|z_j|+\|\mathcal R_{j|j}^{\rm tor}\|\).
On zero-total cell vectors the full reference observation covariance is
\(\Sigma^0_{n,s}=I+aB_sA^+B_s^*\), and
\(v^0_{n,s}=\ip{f_{n,s}}{\Sigma^0_{n,s}f_{n,s}}\) is the reference
variance of \(X_{n,s}\).  Complete the source \(wX_{n,s}\) before the
entry averaging.  It contributes \(e^{v^0_{n,s}w^2/2}\) and translates
the observation argument of the entire interaction by
\(w\Sigma^0_{n,s}f_{n,s}\).  Cell translation and reflection symmetry
make this a first cell cosine.  The corresponding microscopic first
cosine has a nonzero cell average, so choose \(F_{n,s}(w)\) with
\(\alpha B_sF_{n,s}(w)=w\Sigma^0_{n,s}f_{n,s}\).
The finite cosine sum shows that \(F_{n,s}(w)=w\gamma_{n,s}q_n\), with
\(\gamma_{n,s}\) in a fixed compact subinterval of \((0,\infty)\)
when \(n\ge sP\), after \(P\) and the entry threshold are large.
Choose the interval in \eqref{trj:neutral-readout} to contain these
amplitudes.  The exact map identities then give, uniformly on a fixed
source disk,
\begin{equation}\label{trj:gradient-diagnostic}
 \begin{split}
 \log\widehat\E_{j,k} e^{wX_{n_j,s_k}}-\tfrac12v^0_{n_j,s_k}w^2
 &=\mathcal W_j(K_{j|j}^{\rm tor};F_{n_j,s_k}(w))\\
 &=\tfrac12 t_j\ip{F_{n_j,s_k}(w)}{AF_{n_j,s_k}(w)}
            +O(\|\mathcal R_{j|j}^{\rm tor}\|+\nu_j^2).
 \end{split}
\end{equation}
The first equality uses the same extraction scalars in numerator and
denominator, so they cancel.  The estimate includes two source
derivatives.  The uniform mean integral removes the fundamental term
at linear order, and
\(\ip{F_{n,s}(w)}{AF_{n,s}(w)}/w^2\to a\mathfrak g\) as \(s/n\to0\).
This is the physical use of \eqref{trj:neutral-readout}; the finite
cosine estimates, including their errors, are derived in the
overlap comparison below.

There is also a mixed phase diagnostic; write \(s=s_k\) in this paragraph.
On the same activity ball and
source disk, define the normalized phase functional
\begin{equation}\label{trj:phase-functional}
 \mathcal C_j(K;F)=
 \frac{\E_j^{\rm ter}
       [e^{i\alpha\theta}\cZ_{L^j}(K;\phi+F)]}
      {\E_j^{\rm ter}\cZ_{L^j}(K;\phi)}.
\end{equation}
The inserted factor has modulus one, so the same terminal bounds give
joint analyticity and uniform Cauchy bounds on the fixed inner domains.
Changing the common phase by \(\omega/2\) gives
\(\mathcal C_j(\mathcal S K;F)=-\mathcal C_j(K;F)\); the denominator is
invariant.  Let \(\tau_{j,s}\) be the variance of the integrated high
mean in height units, including the mean observation noise.  These
means are independent of the retained mean and of the mean-free source.
The factor \(e^{\tau_{j,s}/2}\) and its inverse are bounded.  For the
actual \(K=K_{j|j}^{\rm tor}\), the completion in
\eqref{trj:gradient-diagnostic} identifies the full physical phase
functional, after removal of the Gaussian prefactor, with
\(e^{\tau_{j,s}/2}\mathcal C_j(K;F_{n_j,s}(w))\).  This formula defines
its analytic, odd extension to the activity ball.  Subtracting the
value at \(w=0\) gives
\begin{equation}\label{trj:mixed-diagnostic}
 \begin{split}
 &e^{-v^0_{n_j,s}w^2/2}\widehat\E_{j,k}
       [\chi_{n_j,s}e^{wX_{n_j,s}}]
       -\widehat\E_{j,k}\chi_{n_j,s}\\
 &\qquad=e^{\tau_{j,s}/2}
       [\mathcal C_j(K;F_{n_j,s}(w))-\mathcal C_j(K;0)]\\
 &\qquad=z_jM_{j,s}(w)+O(\|\mathcal R_{j|j}^{\rm tor}\|+\nu_j^2).
 \end{split}
\end{equation}
The second equality is the terminal expansion, where
\begin{equation}\label{trj:phase-multiplier}
 M_{j,s}(w)=p_{j,s}\operatorname{ave}_x
                  (e^{-i\alpha F_{n_j,s}(w,x)}-1),
 \qquad
 p_{j,s}=\frac{D_0^2}{2}
          e^{\tau_{j,s}/2-aD_{{\rm tail},j}/2}.
\end{equation}
Here \(D_{{\rm tail},j}\) is the diagonal of the remaining mean-free
field covariance.
The average in \eqref{trj:phase-multiplier} is over the microscopic
torus sites.
This is \cite[Equation~(10.4)]{B}.  To see the sign of its high-mean
factor, averaging the interaction over an independent Gaussian constant
multiplies its first Fourier coefficient by \(e^{-\tau_{j,s}/2}\).
The physical mean insertion must undo that averaging, giving
\(e^{\tau_{j,s}/2}\).  Both \(p_{j,s}\) and its inverse are bounded.
For a fixed small real \(w\ne0\), the average in
\eqref{trj:phase-multiplier} stays strictly negative when \(s/n_j\)
is small, so \(|M_{j,s}(w)|\) is bounded below.  In particular both
scalar directions have nonzero diagnostic multipliers.

\begin{lemma}[Finite-history confinement and vanishing]\label{trj:vanishing}
For every sufficiently small prescribed tube, \(R\) and the entry
threshold can be chosen so that, for all sufficiently large \(k\),
every plane state through \(j\le4\cdot2^k\) stays in that tube.
For each \(j\in\mathcal J_k\), the torus states at all stages
\(j_0(k)\le i\le j\) stay in the corresponding small ball.
For every fixed \(0<c<c'<4\), as \(k\to\infty\),
\begin{equation}\label{trj:vanishing-ranges}
 \begin{split}
 \sup_{\substack{i\in\Z,\ j_0(k)\le i\\c2^k\le i\le c'2^k}}
 \bigl(|t_i^{(k)}|+|z_i^{(k)}|
                   +\|\mathcal R_i^{(k)}\|_{i,k}\bigr)&\longrightarrow0,\\
 \sup_{\substack{i\in\Z,\ j\in\mathcal J_k,\ j_0(k)\le i\le j\\
                  c2^k\le i\le c'2^k}}
 \bigl(|t_i^{(k)}|+|z_i^{(k)}|
                   +\|\mathcal R_{i|j}^{\rm tor,(k)}\|\bigr)&\longrightarrow0.
 \end{split}
\end{equation}
\end{lemma}

\begin{proof}
Write \(\varepsilon\) for the entry norm, and call
\(|t_j|+|z_j|+\|\mathcal R_j\|\) the plane full size and
\(|t_i|+|z_i|+\|\mathcal R_{i|j}^{\rm tor}\|\) the full size at stage
\(i\) of a torus stopping at \(j\).  Use the first-exit argument of
\cite[Proposition~10.2]{B} with the
physical limits supplied by Lemma~\ref{trj:physical-torus}.  We give
the finite-history details.  Before a possible first exit from a tube
of radius \(\rho\), the separate remainder on a valid torus satisfies,
by \eqref{trj:triangular} and its nonlinear bound,
\[
 \|\mathcal R_{j|j}^{\rm tor}\|
       \le C\vartheta^{j-j_0(k)}\varepsilon+C\rho^2.
\]
Terms linear in a small remainder have been absorbed in a slightly
larger \(\vartheta<1\).  The forcing uses the same previous plane
singletons, by copying of their full padded neighborhoods.  The plane
remainder has the analogous bound.  At the candidate exit the bounded
map still leaves all new sizes within a fixed multiple of \(\rho\).
Choose \(\rho\) so that this multiple lies inside the readout ball.
The variance derivative of \eqref{trj:gradient-diagnostic} and the
mixed diagnostic \eqref{trj:phase-multiplier}, using their nonzero
multipliers, then bound the new scalars by
\(C(\varepsilon+\rho^2+\delta_{\rm diag})\).  Here the physical
discrepancy \(\delta_{\rm diag}\) can be made arbitrarily small on
all valid stops by taking \(P\), then \(s_k\), large, by
Lemma~\ref{trj:physical-torus} and \eqref{trj:cutoff-error}.
Choose \(\rho\) with \(C\rho\ll1\), and then entry and discrepancy
smaller than a fixed small multiple of \(\rho\).  This excludes the
exit.  The finitely many invalid early tori were handled after
\eqref{trj:valid-stops}.  The same contraction, driven by the now
controlled plane singletons, bounds every stage of each later valid
torus.  Equivalence of the coordinate norms on
the small ball gives the asserted normal-coordinate tube as well.

Here is the uniform vanishing argument, for which no quantitative
physical convergence is needed.  Let \(q_*\) be the supremum, over
fixed \(0<c<c'<4\), of the limsup as \(k\to\infty\) of the largest
plane full size or terminal full size (stage equal to stop) on
\([c2^k,c'2^k]\), with the latter restricted to valid stops.  This is at most
a fixed multiple of \(\rho\).  To estimate remainders on such an
interval, start their contraction at \((c/2)2^k\).  Its homogeneous
term tends to zero exponentially, and the limsup of its forcing is at
most \(Cq_*^2\): after absorbing terms linear in the small torus
remainder, that forcing is bounded by \(C(|t_i|+|z_i|)^2\), using only
the plane singletons.  This also estimates a torus remainder without
identifying it with the plane one.  The two diagnostics have discrepancy
tending uniformly to zero on the interval, since \(n_j\to\infty\)
and \(s_k/n_j\to0\) there.  Their errors and nonzero multipliers
therefore give the same \(Cq_*^2\) bound for the scalar limsups.
Taking the supremum gives \(q_*\le Cq_*^2\).  The tube was chosen below
\(1/C\), hence \(q_*=0\).  For a stage \(i\in[c2^k,c'2^k]\) behind
any valid stop \(j\ge i\), contract that torus's own remainder from
\(\lfloor(c/2)2^k\rfloor\), which lies after the entry for large \(k\).
The homogeneous term tends to zero exponentially, and the forcing uses
the now uniformly vanishing plane singletons on a slightly longer
fractional interval.  The contraction constants are independent of
the stop.  This proves the second supremum in
\eqref{trj:vanishing-ranges}.
\end{proof}

\subsection{The odd signal and comparison of histories}

The phase diagnostic also gives a lower bound, on the absolute scale,
which prevents an arbitrarily late entry from losing the odd direction.
Let \(B_n=\{0,\ldots,n-1\}^2\).  For
\(e\in\{\pm e_1,\pm e_2\}\), let \(\rho_{n,e}\) be the uniform
probability on its face with outward normal \(e\), including the
\(n\) vertices of that face.  Put
\[
 \begin{gathered}
 \bar\rho_n=\frac14\sum_e\rho_{n,e},\qquad
 \rho_n^{\rm vol}=n^{-2}\mathbf 1_{B_n},\\
 z_n^{\rm face}=\E_{B_n}^{h,\rm free}e^{i\ip{h}{\bar\rho_n}},
 \quad
 z_n^{\rm vol}=\E_{B_n}^{h,\rm free}e^{i\ip{h}{\rho_n^{\rm vol}}}.
 \end{gathered}
\]
where \(\E_{B_n}^{h,\rm free}\) is the free Bessel height law at
\(b_c\), modulo its common \(2\pi\) translation.  Both profiles have
total one, and these centered characteristic coefficients are positive.
The critical face estimate \cite[Equation~(2.15)]{C}, using \(a=8\pi\),
says that for every fixed sufficiently large \(G=L^r\),
\[
 \liminf_{n\to\infty}z_{Gn}^{\rm face}/z_n^{\rm face}\ge c_6>0,
\]
where \(c_6\) is independent of \(r\).  Replace it by a smaller number
in \((0,1)\).  For fixed \(r\), iterate the eventual lower bound
\(c_6/2\) on the \(r\) residue classes of the absolute index in
\(n_j=D_0L^j\).  It follows that
\[
 \limsup_{j\to\infty}\frac{-\log z_{n_j}^{\rm face}}{j}
       \le \frac{\log(2/c_6)}{r}.
\]
Letting \(r\to\infty\) proves \(-\log z_{n_j}^{\rm face}=o(j)\).
The difference between the face and uniform profiles has bounded
primary reference cost.  Indeed, for one face route each site's mass
\(n^{-2}\) to its projection on that face along a row or column.
Each of \(O(n^2)\) edges carries at most \(1/n\), so the squared flow
energy is \(O(1)\); averaging the four flows gives the same bound for
\(\bar\rho_n-\rho_n^{\rm vol}\).  The phase-transfer inequality
\cite[Equation~(2.2)]{C} therefore gives
\[
 \bigl|\sqrt{-\log z_n^{\rm face}}
       -\sqrt{-\log z_n^{\rm vol}}\bigr|\le C,
\]
and hence \(-\log z_{n_j}^{\rm vol}=o(j)\).  Restoring the periodic
bonds increases this centered characteristic by \cite[Lemma~2.2]{C}.
We have proved
\begin{equation}\label{trj:phase-lower}
              \E_{n_j}^{\rm tor} e^{i\bar h}\ge e^{-o(j)}.
\end{equation}
This is a statement about one physical torus law at each absolute \(j\).
The noise factor in Lemma~\ref{trj:physical-torus} is bounded below,
and the cutoff error \eqref{trj:cutoff-error} is negligible compared
with \eqref{trj:phase-lower} on any fixed-fraction interval.  Thus
\(|\widehat\E_{j,k}\chi_{n_j,s_k}|\ge e^{-o(j)}\) uniformly there.

\begin{lemma}[Odd cone and transverse remainders]\label{trj:cone}
For every fixed \(0<c<c'<3\), uniformly on
\(c2^k\le j\le c'2^k\) for large \(k\),
\begin{equation}\label{trj:cone-bounds}
 \begin{gathered}
 Y_j=|Z_j|>0,\qquad Y_j\ge e^{-o(j)},\qquad
                       Y_{j+1}/Y_j=1+o(1),\\
 \|Q_{e,j}\|\le C Y_j^2,\qquad
 \|Q_{o,j}\|\le C(|u_j|+Y_j^2)Y_j.
 \end{gathered}
\end{equation}
All little-oh statements are uniform on the indicated intervals.
\end{lemma}

\begin{proof}
Use a slightly longer fixed-fraction interval behind the one in the
statement; Lemma~\ref{trj:vanishing} makes the tube size on it tend
to zero.  Fix a small \(\epsilon>0\).  If
\(\|Q_o\|\le\epsilon|Z|\) and \(Z\ne0\),
\eqref{trj:z-map} gives \(Z_+/Z=1+o(1)\) with positive ratio, while
\eqref{trj:odd-map} gives
\(\|Q_{o,+}\|\le\vartheta\epsilon|Z|+o(|Z|)\).
For a sufficiently small tube the cone is invariant and the sign of
\(Z\) persists.  Outside that cone,
\(|Z|<\epsilon^{-1}\|Q_o\|\), so \eqref{trj:odd-map} contracts
\(\|Q_o\|\) by \(\vartheta+C\delta/\epsilon<1\) when the tube has
radius \(\delta\).  As long as the cone has not been entered this
contracts the whole odd pair geometrically.  The zero pair stays zero.

We justify why the latter alternative cannot last for a fractional
length of a history.  In the torus stopping at \(j\), let
\(r^o_{i|j}\) be the norm of the odd part of its unbalanced remainder
at stage \(i\).  The odd part of \eqref{trj:balance} gives
\(|z_i|\le C(|Z_i|+\|Q_{o,i}\|)\).  The torus map has no linear
forcing from a singleton into its remainder.  By parity each nonlinear
odd forcing contains an odd factor; all terms linear in \(r^o_{i|j}\)
with a small coefficient can be absorbed in contraction.  Hence, in
the small tube,
\begin{equation}\label{trj:torus-odd-convolution}
 r^o_{i+1|j}\le\vartheta' r^o_{i|j}
                         +C(|Z_i|+\|Q_{o,i}\|),
                  \qquad \vartheta'<1.
\end{equation}
If the plane odd pair were outside the cone (or zero) from \(i_0\)
to \(j\), its geometric decay and this convolution would give
\(|z_j|+r^o_{j|j}\le C\vartheta''^{j-i_0}\) for some
\(\vartheta''<1\).  The constants are uniform in the torus stop.
In the terminal integral for \(\chi\), a term with only even factors
vanishes under the half-period shift of the uniform mean.  The fixed
terminal block count, denominator bounded away from zero, and analytic
terminal bounds therefore imply
\[
       |\widehat\E_{j,k}\chi_{n_j,s_k}|
                 \le C(|z_j|+r^o_{j|j}).
\]
Taking \(j-i_0\) to be a fixed positive fraction of \(j\) contradicts
\eqref{trj:phase-lower}.  Thus the cone is entered on an earlier
fixed-fraction interval and persists throughout the interval of the
statement; the zero alternative is excluded.  This is the needed
torus odd-remainder bridge, rather than an identification of the
plane and torus remainders.

Divide \eqref{trj:odd-map} by the positive lower bound for
\(|Z_+|\) in the cone.  It gives a strict contraction for
\(\|Q_o\|/Y\) with a forcing tending uniformly to zero.  Convolution
from the earlier interval proves \(\|Q_o\|/Y=o(1)\), and then
\eqref{trj:z-map} proves the slow step ratio in
\eqref{trj:cone-bounds}.  Choose a fixed \(r_0<1\), as close to one
as needed below.  On the longer interval, eventually
\(r_0Y_i\le Y_{i+1}\le r_0^{-1}Y_i\).
Convolving \eqref{trj:torus-odd-convolution} from an earlier fractional
index, with \(r_0>\vartheta'\), now gives
\(r^o_{j|j}\le CY_j+e^{-c_7j}\).  The preceding terminal bound and
\eqref{trj:phase-lower} imply \(Y_j\ge e^{-o(j)}\): for every fixed
\(\eta>0\) the physical bound eventually exceeds \(e^{-\eta j}\),
and the homogeneous exponential is smaller if \(\eta<c_7\).

Take also \(r_0^3>\vartheta\).  Convolution of
\eqref{trj:even-map} gives \(\|Q_{e,j}\|\le CY_j^2\), since
\(\sum_{d\ge0}\vartheta^d r_0^{-2d}<\infty\) and the homogeneous
exponential is negligible relative to \(Y_j^2\).  Substitution in
\eqref{trj:u-map}, with \(\|Q_o\|=o(Y)\), gives
\begin{equation}\label{trj:u-drift}
                         |u_{i+1}-u_i|\le C Y_i^2.
\end{equation}
For \(i\le j\) in the longer interval this yields
\[
 |u_i|\le |u_j|+C\sum_{h=i}^{j-1}Y_h^2
       \le |u_j|+C r_0^{-2(j-i)}Y_j^2.
\]
Use this and \(Y_i\le r_0^{-(j-i)}Y_j\) in the convolution of
\eqref{trj:odd-map}.  The sums with ratios
\(\vartheta/r_0\) and \(\vartheta/r_0^3\) are finite, giving
\(\|Q_{o,j}\|\le C(|u_j|+Y_j^2)Y_j\).  Its homogeneous term is
negligible relative even to \(Y_j^3\), again by the subexponential
lower bound.  This proves all assertions.
\end{proof}

We next compare two histories through the same physical torus at a
common absolute index.  The following statement quantifies the
overlap of their useful windows.

\begin{lemma}[Comparison on an overlap]\label{trj:overlap}
Fix \(0<c<c'<3\).  Uniformly as \(\min(k_1,k_2)\to\infty\), for
every common integer \(j\) satisfying
\(c2^{k_i}\le j\le c'2^{k_i}\) for \(i=1,2\), put
\((u_i,Y_i)=(u_j^{(k_i)},|Z_j^{(k_i)}|)\).  Then
\begin{equation}\label{trj:overlap-bounds}
 \begin{split}
 |u_1-u_2|&\le C(Y_1^2+Y_2^2),\\
 |Y_1-Y_2|&\le C\sum_{i=1}^2(|u_i|+Y_i)Y_i.
 \end{split}
\end{equation}
The constants may depend on \(c,c'\) and the fixed parameters.
In particular \(Y_1/Y_2\to1\) uniformly under these conditions.
\end{lemma}

\begin{proof}
The window conditions make \(|k_1-k_2|\) bounded in terms of \(c,c'\).
Consequently each cell and entry side is comparable to \(j\), and
\(s/n_j,m/n_j\le e^{-c_8j}\) for either history.  Both stops are valid
for large epochs.  The strict bounds \(0<c<c'<3\) also leave fixed
earlier fractional intervals for the cone and remainder convolutions,
within the larger range of Lemma~\ref{trj:vanishing}.

We first estimate a stopped torus relative to a moving exact Gaussian
curve in its own history.  For the torus with stop \(j\), let
\(\mathcal R^{g,\rm tor}_{i|j}(d)\) be its exact pure Gaussian remainder
at stage \(i\), initialized at this history's side \(m\) with parameter
\(d\).  Its singleton coefficients copy the pure plane ones.  The
proof of \eqref{trj:gaussian-curve}, and the torus contraction, give
\(\|\mathcal R^{g,\rm tor}_{i|j}(d)\|\le C|d|^2\), with uniform analytic
derivatives.  Define
\[
 \Delta^e_{i|j}=(\mathcal R^{\rm tor}_{i|j})_e
                         -\mathcal R^{g,\rm tor}_{i|j}(u_i),
 \qquad
 \Delta^o_{i|j}=(\mathcal R^{\rm tor}_{i|j})_o.
\]
The equality \(T_i=T_i^g(u_i)\) and the even part of
\eqref{trj:balance} imply
\begin{equation}\label{trj:scalar-from-curve}
 |t_i-t_i^g(u_i)|\le C\|Q_{e,i}\|\le CY_i^2,
 \qquad
 z_i=\ell_i^{-1}Z_i+O((|u_i|+Y_i^2)Y_i).
\end{equation}
In the first estimate \(t_i^g\) is the unbalanced pure singleton.
There is no term depending only on \(u_i^2\) in that difference.

We record the torus estimates that accompany
\eqref{trj:scalar-from-curve}:
\begin{equation}\label{trj:moving-torus}
 \|\Delta^e_{j|j}\|\le CY_j^2,
 \qquad
 \|\Delta^o_{j|j}\|\le C(|u_j|+Y_j)Y_j.
\end{equation}
Here is the local comparison proving them.  Subtract the torus map
at \((t_i^g(u_i),0,\mathcal R^{g,\rm tor}_{i|j}(u_i))\) from that at
the actual state.  At the next stage also subtract the change of the
pure remainder from parameter \(u_i\) to \(u_{i+1}\).  Its analytic
derivative and \eqref{trj:u-drift} bound that change by \(CY_i^2\).
The derivative of the torus remainder map at zero has no singleton
row, by \eqref{trj:triangular}.  Its even nonlinear part contains
odd factors in pairs; terms linear in \(\Delta^e\) with a small
coefficient are absorbed in its contraction.  Initially use the
rough convolution \eqref{trj:torus-odd-convolution}, which gives
\(\|\Delta^o_{i|j}\|\le CY_i+e^{-c_9j}\) on the interval in question.
Together with \eqref{trj:scalar-from-curve}, the even map difference
is consequently bounded by
\[
 \|\Delta^e_{i+1|j}\|
       \le\vartheta'\|\Delta^e_{i|j}\|+CY_i^2+e^{-c_9j}.
\]
Convolving from an earlier fractional interval proves its bound in
\eqref{trj:moving-torus}.  For the odd difference, there is no linear
fundamental forcing, and every nonlinear term contains an odd factor.
The remaining even size is bounded by
\(C(|u_i|+Y_i^2)\), using the just proved even bound and the pure
\(O(u_i^2)\) remainder.  The analytic quadratic bound, even if one
bounds an additional odd factor just by \(CY_i\), therefore gives
\[
 \|\Delta^o_{i+1|j}\|
       \le\vartheta'\|\Delta^o_{i|j}\|
                             +C(|u_i|+Y_i)Y_i.
\]
Convolution using \eqref{trj:u-drift} and the slow ratios, exactly as
in the last paragraph of Lemma~\ref{trj:cone}, proves the second bound.
All homogeneous terms are negligible by \(Y_j\ge e^{-o(j)}\).
This comparison uses a different pure parameter at each stage; using
only an \(O(u_i^2)\) bound for the whole pure remainder would not prove
the first bound in \eqref{trj:moving-torus}.

The terminal variance functional obtained from
\eqref{trj:neutral-response} is analytic and invariant under
\(\mathcal S\) on the inner activity ball.  Its first variation in an odd direction at a
purely even state is zero.  Equations~\eqref{trj:scalar-from-curve}
and~\eqref{trj:moving-torus} therefore show that the actual cutoff
variance differs by \(O(Y_j^2)\) from its value on the full exact
Gaussian curve at \(u_j\).  We compute that value exactly, including
the observation prefactor in \eqref{trj:gradient-diagnostic}.
Let \(C_{\ge m}=A^+\mathsf P_m\) be the remaining mean-free covariance at the
start on this torus.  Apply \eqref{trj:pure-response} at \(d=u\) with
the physical field shift \(F_{n_j,s}(w)\) from
\eqref{trj:gradient-diagnostic}.  The logarithmic correction from the
remaining field is
\begin{equation}\label{trj:exact-gaussian-response}
 \frac{u}{2}\ip{A^{1/2}F_{n_j,s}(w)}{
 (I-uA^{1/2}C_{\ge m}A^{1/2})^{-1}A^{1/2}F_{n_j,s}(w)}.
\end{equation}
The resolvent is uniformly bounded for small \(u\), because
\(0\le A^{1/2}C_{\ge m}A^{1/2}\le I\); the determinants cancel.
This is Gaussian completion with the fixed frequency \(\alpha\),
not a change in frequency.

We make explicit the small discretization errors in this computation.
The squared reference norm of the difference between the observation
cosine and the microscopic cosine was bounded by \(C(s/n_j)^2\) in
Lemma~\ref{trj:physical-torus}.  Hence
\(v^0_{n_j,s}=a\mathfrak g+O(s/n_j+n_j^{-2})\).
Cell symmetry makes the observation covariance translate a cell cosine.
Write \(f_Q=f(x_Q/n_j)\), and use the microscopic cosine \(q_n\)
fixed with the grid above.  Since \(\operatorname{ave}_Q f_Q^2=1/2\), symmetry
and the definition of \(v^0_{n_j,s}\) give
\(\Cov_0(y,X_{n_j,s})=2v^0_{n_j,s}(f_Q)_Q\).
Moreover \(B_sq_n=\beta_{n,s}(f_Q)_Q\), where
\(\beta_{n,s}=1+O((s/n_j)^2+n_j^{-2})\), by the finite cosine sum.
Thus \(F_{n_j,s}(w)=2wv^0_{n_j,s}q_n/(\alpha\beta_{n,s})\).
Here \(\ip{q_n}{Aq_n}=2n_j^2\sin^2(\pi/n_j)\) and
\(\mathfrak g=1/(8\pi^2)\).  These formulas give
\[
 \ip{F_{n_j,s}(w)}{AF_{n_j,s}(w)}/w^2
             =a\mathfrak g+O(s/n_j+n_j^{-2}).
\]
Finally the cosine eigenvalue of
\(A^{1/2}C_{\ge m}A^{1/2}=\mathsf P_m\) is
\(1+O((m/n_j)^2)\), directly from the polynomial expansion at zero.
For the overlapping histories all these errors are \(O(e^{-c_8j})\).
Thus the variance predicted by the exact curve is, uniformly for small
\(u\),
\begin{equation}\label{trj:common-variance}
 \mathcal V(u)+O(e^{-c_8j}),\qquad
 \mathcal V(u)=a\mathfrak g+\frac{u a\mathfrak g}{1-u}
                         =\frac{a\mathfrak g}{1-u}.
\end{equation}
Its derivative is bounded away from zero.

Both actual cutoff variances can be compared to the same finite physical
variance \(\Var_n^{\rm tor}x_n\).  The reference-cost sampling estimate
and Cauchy--Schwarz bound the variance difference by \(Cs/n_j\), and
\eqref{trj:cutoff-error} bounds the cutoff change.  Their difference
is therefore \(O(e^{-c_8j})\), without estimating how close that physical
variance is to its limit.  Comparing \eqref{trj:common-variance} for
the two histories, and using the nonzero derivative, gives
\[
 |u_1-u_2|\le C(Y_1^2+Y_2^2)+Ce^{-c_8j}.
\]
The exponential is absorbed by the subexponential lower bounds for
\(Y_1,Y_2\).  This proves the first line of
\eqref{trj:overlap-bounds}.

For the second line use the phase signal at zero tilt.  The same mean
integration that gives \eqref{trj:phase-multiplier} gives its linear
coefficient \(p_{j,s}\) on \(z_j\).  The phase functional
\eqref{trj:phase-functional} is odd, so
its value on the pure curve is zero.  Expanding about that curve,
\eqref{trj:scalar-from-curve}, \eqref{trj:moving-torus}, and the
terminal analytic bounds give
\begin{equation}\label{trj:phase-on-curve}
 \widehat\E_{j,k}\chi_{n_j,s}
       =p_{j,s}\ell_j^{-1} Z_j
                     +O((|u_j|+Y_j)Y_j).
\end{equation}
For example the change of the fundamental derivative from its value
at zero activity is \(O(|u_j|)\), and every other surviving term has
an odd factor.  This accounts for every possible linear odd remainder.

At the same absolute stop, \(D_{{\rm tail},j}\) in
\eqref{trj:phase-multiplier} is identical for the two histories.
The polynomial continuations of the high pieces through side \(L^j\)
telescope to the same zero-mode continuation, independently of \(m\).
The only remaining difference in \(\tau_{j,s}\) is the mean-noise
variance \((s/n_j)^2\).  Consequently
\(p_{j,s_1}-p_{j,s_2}=O(e^{-c_8j})\), while \(p_{j,s_i}\) are bounded
away from zero and \(\ell_j=1+O(e^{-c_1j})\) is common.  The two left
sides of \eqref{trj:phase-on-curve} differ from the same physical
\(\E_n^{\rm tor}e^{i\bar h}\) only by their mean-noise factors and
their cutoff errors, again \(O(e^{-c_8j})\).  Subtracting gives
\[
 |Z_1-Z_2|\le C\sum_{i=1}^2(|u_i|+Y_i)Y_i+Ce^{-c_8j}.
\]
Absorb the exponential as before and use
\(|Y_1-Y_2|\le|Z_1-Z_2|\).  Since \(|u_i|+Y_i\to0\), the resulting
bound is \(|Y_1-Y_2|\le o(1)(Y_1+Y_2)\); it also proves the ratio
assertion.  Both overlap comparisons cancel the unknown finite physical
diagnostic.  They require no convergence rate for the height limit.
\end{proof}

\subsection{The recurrence with dyadic changes of history}

We finish the proof of Proposition~\ref{trj:trajectory}.  At absolute
index \(j\) choose \((u_j,Y_j)\) from \(\mathcal H^{(k)}\) when
\(2^k\le j<2^{k+1}\).  These spliced scalars are positive in the
second coordinate and tend to zero.  At a step that stays in one history,
Lemmas~\ref{trj:parity-interface}, \ref{trj:coefficient-rate}, and
\ref{trj:cone} give
\begin{align}
 u_{j+1}-u_j
   &=-Y_j^2\{H+O(|u_j|+Y_j+(1+j)^{-2})\},
                                      \label{trj:reduced-u}\\
 \frac{Y_{j+1}}{Y_j}
   &=1-b_*u_j+O(u_j^2+Y_j^2+|u_j|(1+j)^{-2}).
                                      \label{trj:reduced-y}
\end{align}
At a crossing \(j+1=2^k\), first take the old history's ordinary step
and then replace it at the same absolute \(j+1\) by the new history.
Lemma~\ref{trj:overlap} adds at most \(C Y_j^2\) to
\eqref{trj:reduced-u}, and at most \(C(|u_j|+Y_j)\) to the ratio in
\eqref{trj:reduced-y}.  Here the slow step ratio and the overlap ratio
make the old and new \(Y\)'s comparable.  These bounds allow the
\(CY_j\) term at a switch before \(u_j\) and \(Y_j\) are known to be
comparable.

We give the barrier argument with these switch errors.  Sample the tail
at fixed lengths \(K\).  At most one switch occurs in a sampled step
when the initial index is sufficiently large.  For fixed \(K\), the
slow ratios give \(Y_i/Y_j=1+o(1)\) throughout the step, and its change
in \(u\) is \(O_K(Y_j^2)\).  Choose \(K\) larger than a fixed multiple
of the constant in the switch error.  Summing
\eqref{trj:reduced-u}, then going sufficiently far into the tail,
gives, with \(+\) now denoting one sampled step,
\begin{equation}\label{trj:sampled-u}
                  -C_2 K Y^2\le u_+-u\le -C_1 K Y^2
\end{equation}
for fixed \(C_1,C_2>0\).  The sampled \(u\)'s decrease to zero,
so they are positive.  Summing the ratio equation and absorbing the
possible \(C|u|\) switch error into its coefficient gives
\begin{equation}\label{trj:sampled-y}
 \frac{Y_+}{Y}=1-\widetilde b u+\mathcal E,
 \qquad C_3K\le\widetilde b\le C_4K,\qquad
 |\mathcal E|\le C_5Y+C_KY^2.
\end{equation}
The constants \(C_1,\ldots,C_5\) can be fixed for all sufficiently
large \(K\); the tail threshold and \(C_K\) may depend on \(K\).
To justify this form, use \(u_i=u+O_K(Y^2)\) inside the step and expand
the product of its \(K\) ratios.  The coefficient errors
\(O(|u_i|(1+i)^{-2})\) sum to
\(O_K(uj^{-2}+Y^2j^{-2})\), since the sampled \(u\) is positive
and \(|u_i|\le u+C_KY^2\).  Absorb the first part into
\(\widetilde b\) and the second into \(\mathcal E\).
The remaining ordinary errors and products are
\(O_K(u^2+uY+Y^2)=O_K(u^2+Y^2)\); their \(u^2\) part changes
\(\widetilde b\) by \(O_K(u)\).  The switch term splits into a part
bounded by \(Cu\), absorbed there as well, and a remainder bounded by
\(C_5Y+C_KY^2\).  Thus
\(\widetilde b=Kb_*+O(1)+O_K(u+j^{-2})\).
Taking \(K\) large and then the tail sufficiently late leaves the
displayed positive coefficient interval.

Choose \(c_->0\) small.  In the region \(0<u\le c_-Y\),
\eqref{trj:sampled-u}--\eqref{trj:sampled-y} imply
\[
 (c_-Y_+-u_+)-(c_-Y-u)
 \ge \{C_1K-C_4Kc_-^2-C_5c_--C_Kc_-Y\}Y^2.
\]
Choose \(c_-\) so the quadratic loss is small, take \(K\) large for
the fixed switch loss, and then take the tail large for the last loss.
The right side is a positive multiple of \(KY^2\).  Thus a nonnegative
lower gap becomes strictly positive and keeps increasing as long as it
is in this region; the same inequality keeps it in the region.  This
contradicts \(u,Y\to0\).  Similarly, in the region \(u\ge c_+Y\),
\[
 (u_+-c_+Y_+)-(u-c_+Y)
 \ge \{-C_2K+C_3Kc_+^2-C_5c_+-C_Kc_+Y\}Y^2.
\]
Choose \(c_+\) sufficiently large, then \(K\) and the tail as above.
The upper gap also increases strictly, again contradicting convergence.
It follows that \(c_-Y<u<c_+Y\) on the sampled tail.  The bounds within
a sampled step give \(u_j\asymp Y_j\) at every sufficiently large index.

Increase the fixed \(K\) once more if needed so that
\(C_3Kc_->2C_5\).  Equations~\eqref{trj:sampled-y} and the two barriers
then give \(Y-Y_+\asymp K Y^2\) on the sampled steps.  Their ratios
tend to one, so taking reciprocals gives
\(1/Y_+-1/Y\asymp K\).  Summing over the sampled indices and filling
the bounded gaps proves
\begin{equation}\label{trj:rough-decay}
                         u_j\asymp Y_j\asymp j^{-1}.
\end{equation}

It remains to read the constant with the required error.  Put
\(c_0=\sqrt{H/b_*}\) and \(I_j=u_j^2-c_0^2Y_j^2\).
At an ordinary step, \eqref{trj:reduced-u}--\eqref{trj:rough-decay}
give
\[
 u_{j+1}-u_j=-HY_j^2+O(j^{-3}),\qquad
 Y_{j+1}-Y_j=-b_*u_jY_j+O(j^{-3}).
\]
In \(I_{j+1}-I_j\), the leading terms
\(-2Hu_jY_j^2+2c_0^2b_*u_jY_j^2\) cancel, and all remaining terms are
\(O(j^{-4})\).  At a switch the extra changes in \(u,Y\) are
\(O(j^{-2})\), so the corresponding bound is \(O(j^{-3})\).
Since \(I_j\to0\), summing its increments to infinity gives
\[
 |I_j|\le C\sum_{i\ge j}i^{-4}
             +C\sum_{2^k\ge j}(2^k)^{-3}
          \le Cj^{-3}.
\]
Together with the positive lower bound for \(u_j+c_0Y_j\) in
\eqref{trj:rough-decay}, this yields
\begin{equation}\label{trj:separatrix-error}
                         u_j=c_0Y_j+O(j^{-2}).
\end{equation}
At ordinary steps the reciprocal form of \eqref{trj:reduced-y} is now
\[
 \frac1{Y_{j+1}}-\frac1{Y_j}=b_*c_0+O(j^{-1}).
\]
At each switch there is an additional \(O(1)\) error.  There are only
\(O(\log j)\) switches up to index \(j\); the ordinary errors also
sum to \(O(\log j)\).  Thus
\(1/Y_j=b_*c_0j+O(\log(2+j))\), and
\begin{equation}\label{trj:u-asymptotic}
                   u_j=\frac1{b_*j}
                       +O\!\left(\frac{\log(2+j)}{j^2}\right).
\end{equation}

For an arbitrary history \(\mathcal H^{(k)}\) at an interior index,
compare it with the spliced history at that same \(j\).  Both indices
are interior in their extended ranges.  The ratio assertion in
Lemma~\ref{trj:overlap} and \eqref{trj:rough-decay} first give
\(Y_j^{(k)}\asymp j^{-1}\); its gradient bound then changes
\eqref{trj:u-asymptotic} by only \(O(j^{-2})\).  Finally,
\eqref{trj:gaussian-curve}, \eqref{trj:scalar-from-curve}, and
Lemma~\ref{trj:cone} give
\begin{align*}
 t_j^{(k)}&=u_j^{(k)}+O(j^{-2}),\qquad z_j^{(k)}=O(j^{-1}),\\
 \|\mathcal R_j^{(k)}\|_{j,k}
  &\le C\{(u_j^{(k)})^2+(Y_j^{(k)})^2+|u_j^{(k)}|Y_j^{(k)}\}
  =O(j^{-2}).
\end{align*}
Substituting \(b_*=2\log L\) proves
Proposition~\ref{trj:trajectory}.  The scale in that conclusion is
the absolute \(j\); no single uniformly valid microscopic start and
no quantitative height scaling limit have been assumed.
